\section{Details of the attackers and of windows}
\label{app:attackers}

% Moved from Sections 4.1, 5.3, 5.4, 6.1 and 6.5 of the main body on 24 October 2026, without changing the
% procedure.

\paragraph{The KeyTAR model.}
KeyTAR reconstructs typed text from keystroke intervals with a fine-tuned language model~\cite{qiu2026keytar}. We use
its public artifact as released (Zenodo record 17254163), with the checkpoint trained on intervals recorded from typing
rather than on intervals recovered through a cache side channel. The artifact is released under the MIT License and the
record under CC~BY~4.0; its statements on ethics and intended use do not exclude reconstruction from the intervals of
public keystroke data. We do not fine-tune it on KLiCKe, and it runs on a local GPU.

\paragraph{A4: KeyTAR.}
A4 takes the intervals between letter, digit and space presses, cuts them into chunks of 20 record segments, gives
each chunk to the public KeyTAR model to reconstruct the text, and chooses the candidate with the smallest normalised
edit distance to the concatenated reconstruction, after both are lowercased and every character other than a letter or
digit is replaced by a space. The model is used with the default settings of its artifact, with intervals capped at
1,023\,ms as in its training data and without preceding context. Among the defaults is a limit of 50 generated tokens
per chunk, which is expected to cut the reconstruction of many chunks short; we kept it as registered rather than
choose, after the fact, a setting that would make the attacker stronger. KeyTAR assumes ordinary English input without
revision~\cite{qiu2026keytar}, it contains a language model, and it does not adapt to the processing. It is therefore
reported in a separate row and is not part of the strongest attacker. The intervals it uses are the same in R0, R1,
R2 and R3, so these forms share one reconstruction.

\paragraph{Alignment and A3.}
Among alignment paths of equal cost, the group listed first is chosen, and where the end can be chosen, the earlier
position in the text. A3 estimates the mean and variance of the logarithm of the press interval for each pair of
consecutive characters; a pair seen fewer than 20 times takes the values of its pair of character types (such as
letter followed by space), and a type pair seen fewer than 20 times takes the values of all intervals. In R3-u the
A1$'$ replay gives the sequence of ``char'' items of the final draft, matched one to one from the start with all
characters of the candidate; the timing term uses the intervals between adjacent final characters that were typed
consecutively, and the length term is the log-likelihood of the difference between the replayed count and the
candidate's character count, estimated on the training logs. The timing term is used both as a sum and as a mean per
term, and the better version is reported as A3 on R3-u; a version with the length term alone separates the
contribution of timing. The check on synthetic logs showed that the one-to-one correspondence in R3-u uses in effect
only the part of the log before the first replay error. In R7, A3 scores the duration of each segment, from its first
to its last remaining press, with a regression on the logarithm of the candidate segment's character count, estimated
on training logs with the same suppression. Details that the registration left open, such as how an interval is
scored when the candidate word has fewer characters than the segment has intervals, were fixed on synthetic logs
before the data were opened and are documented with the code.

\paragraph{The attackers on windows.}
On windows, A0 counts the window's presses as on full texts and compares the count with the candidate window. A1 and
A1$'$ are not run, since a part of the process does not determine the final draft. A2 and A2-b align the window's
segments with the candidate window, with the window end of the alignment; A2 is also run on R0, R1, R5 and R6, whose
key categories allow segments to be built. A3 uses the A2 alignment (in R3-u the A2-b alignment, and the timing term
alone, since the length term is the count of A0). On full texts A3 on R3-u uses the one-to-one correspondence of the
replay, so its difference between full texts and windows also reflects the change of mechanism. A4 reconstructs the
window's intervals and compares the reconstruction with the candidate window.

\paragraph{Candidate windows.}
Windows contain $k = 5$, $10$, $20$ or $50$ non-empty segments (\S\ref{sec:attackers:windows}), with five window starts
per log drawn uniformly from the valid starts. The true window is found once, on the unprocessed log, and used by
every form and attacker. Replaying the log gives, for each press, the character of the final draft that it became; the
true window is the $k$ words of the text that start at the word holding most of the surviving characters of the
window's first segment. Windows whose first segment leaves no character in the final draft, and windows that would run
past the end of the text, are not used. For writers whose rebuilt draft does not match the distributed text, the words
of the two are aligned after normalisation of whitespace, and only windows whose first word maps to a word of the
distributed text are used; the windows that do not map are counted. Each decoy contributes one window of $k$ words,
drawn uniformly from its valid windows. Logs with no valid window of length $k$ are left out for that $k$, and their
number is reported.

\paragraph{Position of true windows.}
To see whether true windows sit at positions that favour the alignment, we show the distribution of the A2 alignment
cost at the true window next to that of the whole log.
